% !TeX document-id = {1c0b4298-276a-4038-8e08-c4a1f4846da7}
% !TeX TXS-program:bibliography = txs:///biber
\documentclass[12pt,a4paper]{article}
\usepackage[T1]{fontenc}
\usepackage{graphicx}
\usepackage{mathtools}
\usepackage{amssymb}
\usepackage{amsthm}
\usepackage{thmtools}
\usepackage{xcolor}
\usepackage{nameref}
\usepackage{hyperref}
\usepackage{color}
\usepackage{float}
%\usepackage[margin=2.5cm]{geometry}
\usepackage[brazilian]{babel}
\usepackage[
backend=biber,
style=authoryear,
natbib=true
]{biblatex}
\addbibresource{../bibliography.bib}
\newtheorem{definition}{Definição}
\newtheorem{lemma}{Lema}
\newtheorem{proposition}{Proposição}
\newtheorem{fact}{Fato}
\title{Método Delta}
\author{ Professor Luís Antonio Fantozzi Alvarez}
%\date{21 de fevereiro de 2025}
\begin{document}
	\maketitle
	
	Nestas notas, vamos demonstrar a versão do método Delta vista em aula:
	
		\begin{lemma}[Método Delta]
		Seja $\boldsymbol{X}_n$ uma sequência de vetores aleatórios com contradomínio em $\mathbb{R}^k$. Seja $\theta \in \mathbb{R}^k$, e $(r_n)_n \in \mathbb{R}^\infty$ uma sequência de números reais positivas tais que $r_n \to \infty$. Suponha que $r_n(\boldsymbol{X}_n - \theta)\overset{d}{\to} S$, onde $S$ é um vetor aleatório com valores em $\mathbb{R}^k$. Se $\psi: \mathbb{R}^k \mapsto \mathbb{R}^l$ é continuamente diferenciável numa vizinhança de $\theta$, então:
		$$r_n(\psi(\boldsymbol{X}_n)-\psi(\theta)) \overset{d}{\to} \nabla_{x}\psi(\theta) S\, ,$$
		onde $\nabla_{x}\psi(\theta)$ é o Jacobiano de $\psi$ avaliado em $\theta$.
	\end{lemma}
	\begin{proof}
		Como $\psi$ é continuamente diferenciável numa vizinhança de $\theta$, sabemos que existe $\epsilon > 0$ tal que $s \mapsto \psi_x(s)$ existe e é contínua no conjunto $B(\theta;\epsilon) := \{x \in \mathbb{R}^k: \lVert x-\theta \rVert < \epsilon \}$. Considere a sequência de eventos $E_n = \{\boldsymbol{X}_n \in B(\theta;\epsilon)\}$, $n \in \mathbb{N}$. Observe que, trivialmente, podemos escrever:
		
		$$r_n(\psi(\boldsymbol{X}_n)-\psi(\theta)) = \mathbf{1}_{E_n}r_n(\psi(\boldsymbol{X}_n)-\psi(\theta)) + \underbrace{\mathbf{1}_{E_n^\complement} r_n(\psi(\boldsymbol{X}_n)-\psi(\theta))}_{=: R_n}\, .$$
		
		Note, ademais que, como $\boldsymbol{X}_n \overset{p}{\to} \theta$ (vimos isso em aula), para todo $\xi > 0$:
		
		$$0 \leq \mathbb{P}[\{\lVert R_n \rVert>\xi\}] \leq \mathbb{P}[E^\complement_n] \to 0 \implies \lim_{n \to \infty} \mathbb{P}[\{\lVert R_n \rVert>\xi\}] = 0\, ,$$
		de onde concluímos que $ R_n \overset{p}{\to} \boldsymbol{0}_{l \times 1}$. Em seguida, observamos que, como $\psi$ é continuamente diferenciável em $B(\theta;\epsilon)$, para todo $i =1\, \ldots, l$, temos que, pelo teorema do valor médio:\footnote{Para os nossos fins, é suficiente a seguinte versão: seja $O \subseteq \mathbb{R}^k$ um subconjunto aberto e convexo, e $f: O \mapsto \mathbb{R}$ uma função continuamente diferenciável (i.e. cujas derivadas parciais existem e são contínuas). Então, para todo $a, b \in O$, existe $\lambda \in [0,1]$ tal que, definindo $c = \lambda a + (1-\lambda)b$:
		
		$$f(a)- f(b) = \nabla_x f(c) (a-b)\, .$$
		 }
		
		$$\mathbf{1}_{E_n}r_n(\psi_i(\boldsymbol{X}_n)-\psi_i(\theta)) = \mathbf{1}_{E_n}\nabla_x \psi_i(\tilde{\theta}_{i,n})r_n(\boldsymbol{X}_n - \theta)\, ,$$ 
		onde $\tilde{\theta}_{i,n}$ é uma variável aleatória que, com probabilidade 1, se encontra no segmento de reta que une $\boldsymbol{X}_n$ a $\theta$. Empilhando a relação acima, obtemos:
		
			$$\mathbf{1}_{E_n}r_n(\psi(\boldsymbol{X}_n)-\psi(\theta)) = \mathbf{1}_{E_n}r_n \begin{pmatrix}
				\psi_1(\boldsymbol{X}_n) \\
				\psi_2(\boldsymbol{X}_n)\\
				\vdots \\
				\psi_l(\boldsymbol{X}_n)
			\end{pmatrix} = \mathbf{1}_{E_n}  \begin{pmatrix}
			\nabla_x \psi_1(\tilde{\theta}_{1,n}) \\
		\nabla_x \psi_2(\tilde{\theta}_{2,n})\\
			\vdots \\
		\nabla_x \psi_l(\tilde{\theta}_{l,n})
			\end{pmatrix}r_n(\boldsymbol{X}_n - \theta)\, ,$$ 
			
				Agora, pelo teorema do mapa contínuo, temos que, quando $n \to\infty$:
			
			$$\begin{pmatrix}
				\nabla_x \psi_1(\tilde{\theta}_{1,n}) \\
				\nabla_x \psi_2(\tilde{\theta}_{2,n})\\
				\vdots \\
				\nabla_x \psi_l(\tilde{\theta}_{l,n})
			\end{pmatrix} \overset{p}{\to} \begin{pmatrix}
				\nabla_x \psi_1(\theta) \\
				\nabla_x \psi_2(\theta)\\
				\vdots \\
				\nabla_x \psi_l(\theta)
			\end{pmatrix} = \nabla_x \psi(\theta)$$
			Ademais, observe que:
			
			$$\mathbb{E}[|\mathbf{1}_{E_n}-1|] =\mathbb{E}[\mathbf{1}_{E_n^\complement}] =  \mathbb{P}[E^\complement_n] \to 0\, , $$
			isto é, $\mathbf{1}_{E_n} \overset{L_1}{\to}1$, e, como consequência, $\mathbf{1}_{E_n} \overset{p}{\to}1$. Aplicando o lema de Slutsky, obtemos, dessa forma, que 
		
		
			$$\mathbf{1}_{E_n}r_n(\psi(\boldsymbol{X}_n)-\psi(\theta))\overset{d}{\to} \nabla_x \psi(\theta)S,\, $$
			e, como $R_n \overset{p}{\to} \boldsymbol{0}_{l\times 1}$, uma segunda aplicação do lema de Slutsky nos leva a concluir que: 
			
			$$r_n(\psi(\boldsymbol{X}_n)-\psi(\theta))\overset{d}{\to} \nabla_x \psi(\theta)S\, ,$$
			como desejado.
	\end{proof}

\end{document}